\setchoices{label=(\Alph*)}

\begin{question}[index=1]
当 $x\to0$ 时，$\left(\cos\dfrac{x}{3}\right)^3\left(\cos\dfrac{x}{5}\right)^5-\cos x$ 是 $x^2$ 的
\begin{choices}[columns=2]
\item 高阶无穷小
\item 低阶无穷小
\item 同阶无穷小但不等价
\item 等价无穷小
\end{choices}
\end{question}
\QuestionBreak

\begin{question}[index=2]
设 $f\!\left(x+2y,\dfrac yx\right)=\sin x\cdot\ln y$，则 $\left.\dfrac23\dfrac{\partial f}{\partial x}\right|_{(1,1)}+\left.\dfrac{\partial f}{\partial y}\right|_{(1,1)}=$
\begin{choices}[columns=4]
\item $\dfrac13\sin\dfrac13$
\item $\dfrac12\sin\dfrac13$
\item $\dfrac23\sin\dfrac13$
\item $\sin\dfrac13$
\end{choices}
\end{question}
\QuestionBreak

\begin{question}[index=3]
设 $L\colon x^2+y^2=a^2\;(a>0)$，取逆时针方向，$f(x,y)$ 具有二阶连续偏导数，且 $\dfrac{\partial^2 f}{\partial x^2}+\dfrac{\partial^2 f}{\partial y^2}=\eu^{-x^2-y^2}$，则 $\displaystyle\oint_L(x^2+y^2)\left(\frac{\partial f}{\partial x}\dd y-\frac{\partial f}{\partial y}\dd x\right)=$
\begin{choices}[columns=4]
\item $\pi a^2$
\item $a^2\pi(1-\eu^{-a^2})$
\item $a^2\pi(\eu^{-a^2}-1)$
\item $a^2\eu^{-a^2}\pi$
\end{choices}
\end{question}
\QuestionBreak

\begin{question}[index=4]
设 $f(x)$ 在闭区间 $[0,1]$ 上可导，$f(0)=0$ 且 $f'(x)>0$，反常积分 $\displaystyle\int_0^1\frac{f(x)}{x^\alpha}\dd x$ 和 $\displaystyle\int_1^{+\infty}\frac{\ln x}{x^\alpha}\dd x$ 均收敛，则 $\alpha$ 的取值范围为
\begin{choices}[columns=4]
\item $1<\alpha\le2$
\item $1<\alpha<2$
\item $0\le\alpha<2$
\item $0<\alpha<2$
\end{choices}
\end{question}
\QuestionBreak

\begin{question}[index=5]
设 $A$ 为 $4\times5$ 阶矩阵，$B$ 为 $4$ 阶矩阵，若齐次线性方程组 $AA^{\mathrm T}x=0$ 仅有零解，$B^*x=0$ 解向量组的秩为 $3$，则齐次线性方程组 $A^{\mathrm T}Bx=0$ 解向量组的秩为
\begin{choices}[columns=4]
\item $4$
\item $3$
\item $2$
\item $1$
\end{choices}
\end{question}
\QuestionBreak

\begin{question}[index=6]
设 $P,Q$ 为 $3$ 阶矩阵，$P^*$ 为 $P$ 的伴随矩阵，$|P|=2$，$|P-2Q|=6$，则 $|P^*Q-E|=$
\begin{choices}[columns=4]
\item $-3$
\item $-4$
\item $8$
\item $12$
\end{choices}
\end{question}
\QuestionBreak

\begin{question}[index=7]
设四阶矩阵 $A=(a_{ij})_{4\times4}$ 满足 $\displaystyle\sum_{j=1}^4a_{ij}=0$（$i=1,2,3,4$），$M_{ij}$ 与 $A_{ij}$ 分别为元素 $a_{ij}$ 的余子式和代数余子式，下列四个式子中必成立的有多少个？
\par\circlednumber{1}\quad $M_{21}+M_{22}+M_{23}+M_{24}=0$
\par\circlednumber{2}\quad $A_{21}+A_{22}+A_{23}+A_{24}=0$
\par\circlednumber{3}\quad $M_{21}=M_{22}=M_{23}=M_{24}$
\par\circlednumber{4}\quad $A_{21}=A_{22}=A_{23}=A_{24}$
\begin{choices}[columns=4]
\item $1$
\item $2$
\item $3$
\item $4$
\end{choices}
\end{question}
\QuestionBreak

\begin{question}[index=8]
设随机事件 $A,B,C$，$P(BC)>0$ 且 $P(B)<1$，则下列 $3$ 个命题正确的个数为
\par\circlednumber{1}\quad 若 $P(A\mid B)>P(A)$，$P(B\mid C)>P(B)$，则必有 $P(A\mid C)>P(A)$
\par\circlednumber{2}\quad 若 $P(A\mid B)>P(A)$，$P(A\mid C)>P(A)$，则必有 $P(A\mid BC)>P(A)$
\par\circlednumber{3}\quad 若 $P(A)>P(A\mid B)$，则必有 $P(A)<P(A\mid\overline B)$
\begin{choices}[columns=4]
\item $0$
\item $1$
\item $2$
\item $3$
\end{choices}
\end{question}
\QuestionBreak

\begin{question}[index=9]
设总体 $X\dist B(1,p)$，$X_1,\cdots,X_{1000}$ 是来自总体 $X$ 的简单随机样本，$\overline X$ 为样本均值，则下列结论不正确的是
\begin{choices}[columns=2]
\item $\displaystyle\sum_{i=1}^{1000}X_i\dist B(1000,p)$
\item $\displaystyle\sum_{i=1}^{1000}X_i\overset{\text{近似}}{\dist}N(1000p,1000p(1-p))$
\item $\overline X\overset{\text{近似}}{\dist}N(p,p(1-p))$
\item $\displaystyle\lim_{n\to\infty}\overline X\overset{P}{=}p$
\end{choices}
\end{question}
\QuestionBreak

\begin{question}[index=10]
设随机变量 $X$ 的概率密度函数为 $f(x)$，则随机变量 $Y=|X|$ 的概率密度函数为
\begin{choices}[columns=2]
\item $f_Y(y)=\begin{cases}\dfrac{f(y)-f(-y)}2,&y>0,\\0,&y\le0,\end{cases}$
\item $f_Y(y)=\begin{cases}f(y)-f(-y),&y>0,\\0,&y\le0,\end{cases}$
\item $f_Y(y)=\begin{cases}\dfrac{f(y)+f(-y)}2,&y>0,\\0,&y\le0,\end{cases}$
\item $f_Y(y)=\begin{cases}f(y)+f(-y),&y>0,\\0,&y\le0.\end{cases}$
\end{choices}
\end{question}
